\section{本轮新结果（第 15--17 轮）：Stiefel 刚性、能量分解与全局二分}\label{sec:r1517}

本节记录第 15 至第 17 轮的工作。三轮之间有一条清晰的主线：
第~\ref{sec:harm} 节把问题翻译成了「帧测度 $+$ 方向」的语言，并得到了一个\emph{条件性}的六原子归约
（定理~\ref{thm:atomreduce}）。本轮做三件事：
第一，把条件性归约所依赖的临界点方程组\textbf{完整解出}（\S\ref{sec:r17:erratum}--\S\ref{sec:r17:kkt}），
其间修正了一处会导致整套刚性论证被误判为失效的\textbf{代码错误}；
第二，找到主不等式的一个\textbf{最简等价形式}（\S\ref{sec:r17:tau}），它把「$\rho\le\rho_0$」压缩成
一个单一的标量关系 $\tau=\beta-(1-\alpha)$，并配上一条精确的 Frobenius 恒等式（\S\ref{sec:r17:frob}）；
第三，也是本轮的主定理：给出一个\textbf{无条件的全局二分定理}（定理~\ref{thm:dichotomy}），
它把「$N\ge7$ 能否成为极大元」这件事压缩成两个各自良定的子问题，
并且在数值上把这两个分支都\textbf{精确}地实现出来了（\S\ref{sec:r17:num}）。

同时本轮有一条\textbf{否证}：上一轮准备的「第四不等式」候选 $\alpha+\sqrt{A_1A_2}\le\rho_0$
虽然在冠军处尖锐取等，但\textbf{不是全局真命题}（命题~\ref{prop:fourthfalse}）。
按既定规则，否证与正定理同等记录。

\subsection{勘误：Stiefel 梯度公式的一次误报}\label{sec:r17:erratum}

第 15 轮曾报告「解析 Stiefel 梯度与数值差分相差 $2.03$ 倍（随机构型最大相对误差 $10$）」，
并据此把整条 $h_iW_i=\rho\Lam u_i$ 的刚性论证标为可疑。本轮查明：\textbf{这是脚本错误，不是数学错误}。

\begin{proposition}[正部的环境梯度]\label{prop:gradformula}
固定 $g>0$，令 $Q=VV^\top$、$V=U\operatorname{diag}(\sqrt g)$、$Q_+=K\circ Q$，
其中 $K_{ij}=\mathbf 1\{\langle u_i,u_j\rangle>0\}$（$K_{ii}=1$）。
设 $\lambda_{\max}(Q_+)$ 单重、Perron 向量 $x>0$、$\|x\|=1$，并记 $h_i=x_i/\sqrt{g_i}$、$c_i=g_ih_i=\sqrt{g_i}\,x_i$。
只要没有 $\langle u_i,u_j\rangle=0$，$\rho(g,U)=\lambda_{\max}(Q_+)$ 在 $(u_i)$ 方向可微，且
\begin{equation}\label{eq:gradu}
\frac{\partial\rho}{\partial u_i}=2c_iW_i,\qquad
W_i:=\sum_{j\in A(i)}c_ju_j,\qquad A(i):=\{j:\langle u_i,u_j\rangle>0\}\ \ (\text{含 }j=i).
\end{equation}
\end{proposition}

\begin{proof}
单重特征值的包络定理给 $\mathrm d\rho=x^\top(\mathrm dQ_+)\,x$。
在符号模式不变处 $\mathrm d(Q_+)_{ij}=K_{ij}\,\mathrm dQ_{ij}$，故
$\mathrm d\rho=\sum_{ij}x_ix_jK_{ij}\,\mathrm dQ_{ij}=2\bigl\langle (xx^\top\!\circ K)V,\mathrm dV\bigr\rangle$。
而 $\bigl((xx^\top\!\circ K)V\bigr)_i=x_i\sum_{j\in A(i)}x_j\sqrt{g_j}\,u_j=x_iW_i$，
且 $\mathrm dv_i=\sqrt{g_i}\,\mathrm du_i$，代入即得 $\mathrm d\rho=2\sum_ic_i\langle W_i,\mathrm du_i\rangle$。
\end{proof}

\textbf{错误的根源}。脚本用 \texttt{np.maximum(K,0)} 去构造 $W_i$，
即把\emph{截断余弦值} $(\langle u_i,u_j\rangle)_+$ 当成了\emph{指示函数} $\mathbf 1\{\langle u_i,u_j\rangle>0\}$。
这两个矩阵\textbf{只在再与 $K$ 缩并一次时才相同}：
$\sum_j(\langle u_i,u_j\rangle)_+c_j=\sum_{j\in A(i)}c_j\langle u_i,u_j\rangle=\langle u_i,W_i\rangle=p_i$，
而 $\sum_j(\langle u_i,u_j\rangle)_+c_ju_j\ne W_i$。
第~\ref{sec:harm} 节里所有用到 $p_i$ 的恒等式因此\textbf{不受影响}；受影响的只有把 $W_i$ 当向量使用的三处
（梯度、KKT 拟合、椭球面恒等式）。

\begin{center}\small
\begin{tabular}{lcc}
\hline
公式 & 冠军处 $\max_i\lvert$数值$-$解析$\rvert$ & 随机构型 \\
\hline
\eqref{eq:gradu}（指示掩码，正确） & $4.65\times10^{-10}$ & 全部一致 \\
截断余弦当掩码（错误） & $1.87\times10^{-1}$ & 相对误差最大 $10$ \\
\hline
\end{tabular}
\end{center}

\noindent 数值列是环境坐标中心差分（步长 $10^{-6}$）与 \eqref{eq:gradu} 的逐分量比较；
$4.65\times10^{-10}$ 正是该步长的截断误差量级。\textbf{结论：命题~\ref{prop:gradformula} 成立，
第 15 轮的刚性结果全部有效。}

\subsection{Stiefel KKT、椭球面恒等式与枢纽刚性}\label{sec:r17:kkt}

把 \eqref{eq:gradu} 投影到 Stiefel 流形 $\{V:V^\top V=I_3\}$ 的法空间 $\{VS:S=S^\top\}$ 上，
就得到临界点的完整方程组。

\begin{theorem}[Stiefel KKT 与迹一刚性]\label{thm:stiefelkkt}
设 $(U,g)$ 是 $\rho$ 在帧测度簇 $\mathcal M$（定义~\ref{def:framemeas}）上的临界点，
$\lambda_{\max}(Q_+)$ 单重且符号模式非退化。令 $S:=2V^\top(xx^\top\!\circ K)V$。则
\begin{equation}\label{eq:kkt}
2h_iW_i=Su_i\quad(\forall i),\qquad \tr S=2\rho .
\end{equation}
特别地，取 $\Lam:=S/(2\rho)$ 有
\begin{equation}\label{eq:trace1}
h_iW_i=\rho\Lam u_i,\qquad \tr\Lam=1,\qquad h_i^2=u_i^\top\!\Lam u_i,\qquad
\langle u_i,W_i\rangle=p_i=\rho h_i .
\end{equation}
\end{theorem}

\begin{proof}
法空间投影：临界 $\Leftrightarrow$ 存在对称 $S$ 使 $\nabla_V\rho=VS$（逐行 $\nabla_{v_i}\rho=(Sv_i)^\top$）。
由命题~\ref{prop:gradformula} 的证明，$\nabla_V\rho=2(xx^\top\!\circ K)V$，故 \eqref{eq:kkt} 的前一式；
左乘 $V^\top$ 并用 $V^\top V=I_3$ 得 $S=2V^\top(xx^\top\!\circ K)V$。
迹：$\tr S=2\sum_i v_i^\top(x_iW_i)=2\sum_ix_i\sqrt{g_i}\langle u_i,W_i\rangle=2\sum_ic_ip_i=2x^\top Q_+x=2\rho$。
把 $2h_iW_i=Su_i=2\rho\Lam u_i$ 与 $u_i$ 作内积，并用 $\langle u_i,W_i\rangle=p_i$，
得 $2h_ip_i=2\rho\,u_i^\top\!\Lam u_i$；代入 $p_i=\rho h_i$（由 \eqref{eq:kkt} 与 $\tr\Lam=1$ 亦可独立导出）
即得 $h_i^2=u_i^\top\!\Lam u_i$。
\end{proof}

\begin{remark}[迹一归一化的几何意义]\label{rmk:trace1}
$h_i^2=u_i^\top\!\Lam u_i$ 与 $\tr\Lam=1$ 合起来说：
\textbf{Perron 向量的「相对振幅」$h_i$ 恰好是原子方向 $u_i$ 在一个固定迹一椭球 $\Lam$ 上的支撑函数}。
冠军处 $\Lam$ 的谱是 $\{0.594608575866,\ 0.202695712067^{[2]},\ \}$，
即一个\emph{旋转椭球}，长半轴 $1/\sqrt{0.2027}$ 沿枢纽方向；
$\min_{S^2}u^\top\!\Lam u=0.2027>0$，故松弛「$\theta\ge0$ 在整个 $S^2$ 上」在冠军处是\textbf{紧}的。
\end{remark}

\begin{corollary}[椭球面恒等式]\label{cor:ellipsoid}
记 $m:=\sum_ic_iu_i=V^\top x$（于是 $\alpha=|m|^2$），
$S_i:=\sum_{j\in O(i)}c_ju_j$（钝部部分和，$O(i)=\{j:\langle u_i,u_j\rangle<0\}$），
则 $W_i=m-S_i$，且在定理~\ref{thm:stiefelkkt} 的条件下
\begin{equation}\label{eq:ellipsoid}
(m-S_i)^\top\Lam^{-1}(m-S_i)=\rho^2\qquad\text{对一切 }i\text{ 成立},
\end{equation}
等价地 $S_i^\top\Lam^{-1}S_i-2\langle z,S_i\rangle=-D$ 与 $i$ 无关，其中 $z=\Lam^{-1}m$、$D=m^\top\Lam^{-1}m-\rho^2$。
\end{corollary}

\begin{proof}
由 \eqref{eq:trace1}，$u_i=(h_i/\rho)\Lam^{-1}(m-S_i)$；代入 $h_i^2=u_i^\top\!\Lam u_i$ 得
$h_i^2=(h_i^2/\rho^2)(m-S_i)^\top\Lam^{-1}(m-S_i)$，约去 $h_i^2>0$ 即得。
\end{proof}

\begin{remark}[几何读法]\label{rmk:ellipsoidread}
\eqref{eq:ellipsoid} 说：\textbf{所有钝部部分和 $S_i$ 落在以 $m$ 为中心、以 $\Lam^{-1}$ 定形、
半径 $\rho$ 的同一张椭球面上}。这是一个把「$N$ 个原子的组合信息」压缩成「一张二次曲面」的强约束：
$N$ 越大，落在同一张椭球面上的点越多。这是 \S\ref{sec:r17:dichotomy} 全局二分的几何原型。
\end{remark}

\begin{corollary}[枢纽无关式与精确 Cauchy--Schwarz 缺陷]\label{cor:hubfree}
在定理~\ref{thm:stiefelkkt} 的条件下：
\begin{enumerate}\itemsep1pt
\item 若存在\textbf{枢纽}（某 $i_0$ 使 $O(i_0)=\varnothing$，即 $S_{i_0}=0$），则
\begin{equation}\label{eq:hubfree}
\rho^2=m^\top\Lam^{-1}m,\qquad
\tr(\Lam^{-1})-\rho^2=\sum_ig_i\bigl|\Lam^{-1/2}(u_i-h_im)\bigr|^2\ \ge0 ;
\end{equation}
\item $m=(\rho/h_{i_0})\Lam u_{i_0}$，从而 $u_{i_0}$ 是 $\Lam$ 的特征向量、$\lambda_{\max}(\Lam)=h_{i_0}^2$，且
$\alpha=\rho^2h_{i_0}^2$。
\end{enumerate}
\end{corollary}

\begin{proof}
(1) $S_{i_0}=0$ 代入 \eqref{eq:ellipsoid}。第二式是把 $\tr(\Lam^{-1})=\sum_ig_i\,u_i^\top\Lam^{-1}u_i$
（帧条件 $\sum g_iu_iu_i^\top=I_3$）展开并配方。
(2) 由 \eqref{eq:trace1} 在 $i=i_0$：$h_{i_0}m=\rho\Lam u_{i_0}$。
再取 $u_{i_0}$ 与两边内积：$h_{i_0}\langle u_{i_0},m\rangle=\rho h_{i_0}^2$，
即 $\langle u_{i_0},m\rangle=\rho h_{i_0}=p_{i_0}$；结合 $\Lam u_{i_0}=(h_{i_0}/\rho)m$ 知
$m$ 与 $u_{i_0}$ 平行且 $\Lam u_{i_0}=h_{i_0}^2u_{i_0}$，于是 $\alpha=|m|^2=\rho^2h_{i_0}^2$。
\end{proof}

冠军处的数值（全部到机器精度）：
\begin{center}\small
\begin{tabular}{ll}
\hline
检验 & 残差 \\
\hline
$h_iW_i=\rho\Lam u_i$（\eqref{eq:trace1}） & $4.44\times10^{-16}$ \\
$\tr\Lam-1$ & $0$（$\tr\Lam=1.000000000000000$） \\
$h_i^2=u_i^\top\!\Lam u_i$（无常数项，$\lambda_0=0$） & $2.2\times10^{-16}$ \\
$\rho^2-m^\top\Lam^{-1}m$ & $6.66\times10^{-16}$ \\
$\tr(\Lam^{-1})-\rho^2-\sum g_i|\Lam^{-1/2}(u_i-h_im)|^2$ & $<10^{-11}$ \\
$m-(\rho/h_0)\Lam u_0$ & $4.45\times10^{-16}$ \\
$\kappa_{\rm ang}-1$（$u_0$ 是 $\Lam$ 特征向量） & $0$ \\
$\alpha-\rho^2h_0^2$ & $<10^{-12}$ \\
$m$ 与 $u_0$ 的夹角余弦 & $1.000000000000000$ \\
\hline
\end{tabular}
\end{center}

